\documentclass[11pt,a4paper]{article}

\usepackage[utf8]{inputenc}
\usepackage[T1]{fontenc}
\usepackage{lmodern}
\usepackage{amsmath,amssymb}
\usepackage{booktabs}
\usepackage{graphicx}
\usepackage{caption}
\usepackage[margin=1in]{geometry}
\usepackage{microtype}
\usepackage[hidelinks]{hyperref}
\usepackage{url}

\title{\bfseries Counting fills misrepresents liquidations:\\
evidence from a complete on-chain derivatives ledger}

\author{monproweb\\ \small Quasareum\\ \small \texttt{contact@quasareum.com}}

\date{27 July 2026}

\begin{document}
\maketitle

\begin{abstract}
\noindent
Forced liquidations in cryptocurrency derivatives are almost always counted from
exchange-published event feeds, in which one economic liquidation appears as several records.
We quantify the resulting distortion on Hyperliquid, an on-chain perpetual-futures venue whose
node-fills archive is, to our knowledge, the only public liquidation record that is
simultaneously complete and attributed to the liquidated account. On a uniformly stratified
sample of 351{,}540 liquidation episodes reconstructed from 2{,}010{,}042 fills over one year,
counting fills rather than episodes inflates the event count by a factor of \textbf{5.72}. The
inflation is not a constant: it is \emph{size-dependent}, rising from a median of 2 fills per
episode below the median size to a median of 72 in the top percentile, so the bias is
concentrated on precisely the events that liquidation studies are about. The top 1\% of episodes
generate 23.1\% of all fills. We further characterise the episode-size distribution as far as
the data permits, and report a negative result that we believe is the more useful one: the tail
is unambiguously heavy---an exponential is rejected decisively---and it is not a power law, but
it \emph{cannot be named} beyond that. Lognormal and Weibull specifications both dominate a
Pareto fit at every estimable threshold across nearly three orders of magnitude of cut-off, yet
which of the two wins reverses with the threshold, so any named family would report the
analyst's cut-off rather than the distribution. Finally, we document that the fill-counting bias
does \emph{not} transfer to centralised-exchange feeds, which publish orders rather than fills,
but that two of the three major feeds carry a different and, for counting purposes, worse
defect: they publish per-second maxima rather than samples. All data, code and pre-registrations
are public.
\end{abstract}

\section{Introduction}

Empirical work on cryptocurrency derivatives routinely counts liquidation events: how many
occurred, how large they were, how they cluster in cascades, how they propagate across venues.
The counting is nearly always done from an exchange-published liquidation feed, and the unit of
that feed is not the unit of the question.

A trader whose position is force-closed experiences one liquidation. The exchange executes it as
a sequence of fills against the order book. If the published record is per fill, then a single
economic event enters the dataset several times, at several sizes, none of them the size of the
position that was closed. Any statistic computed over those records---a count, a size
distribution, a tail index, an inter-arrival time---measures the exchange's execution mechanics
as much as the market.

This is a well-understood problem in principle. What has been missing is a public dataset in
which the correct unit can be reconstructed, so that the distortion can be measured rather than
assumed. That requires a liquidation record that is (i) complete, with no sampling or rate limit,
and (ii) attributed, so that fills belonging to one account's forced close can be grouped. To our
knowledge no centralised venue publishes both.

Hyperliquid does. It is an on-chain perpetual-futures exchange whose node data includes every
fill, each carrying the liquidated account address, the mark price at liquidation, the starting
position and the realised loss. This paper uses that archive to measure how far fill-counting
departs from episode-counting, and what the episode-size distribution actually looks like.

\paragraph{Contributions.}
\begin{enumerate}
\item We measure a \textbf{5.72$\times$} count inflation from fill-counting on a uniformly
      stratified one-year sample, and show the factor is robust to how ``one liquidation'' is
      defined (Section~\ref{sec:unit}).
\item We show the inflation is \textbf{size-dependent}, growing from 2 to 72 fills per episode
      between the lower half and the top percentile of the size distribution, so it is worst
      exactly where liquidation research concentrates (Section~\ref{sec:count}).
\item We show fill-counting \textbf{compresses the size distribution}, by factors of
      3.4$\times$ at the 99th percentile and 10.9$\times$ at the 99.9th
      (Section~\ref{sec:compression}).
\item We characterise the episode-size tail and report that it \textbf{cannot be named}:
      heavy, not a power law, but with the ranking between lognormal and Weibull reversing
      across the threshold range (Section~\ref{sec:tail}).
\item We document what centralised-exchange feeds actually publish, and show the failure mode
      there is different and worse for counting purposes (Section~\ref{sec:cex}).
\end{enumerate}

\paragraph{On the negative result.} Section~\ref{sec:tail} sets out to name the tail and fails.
We report the failure in full because the alternative---quoting whichever family wins at the
threshold one happened to choose---is how a threshold choice becomes a stylised fact. Three
tail-index estimates were produced during this work and all three were withdrawn; the record of
those withdrawals is public and is described in Section~\ref{sec:repro}.

\section{Setting and data}
\label{sec:data}

\subsection{Hyperliquid liquidations}

Hyperliquid is a perpetual-futures venue with an on-chain order book. Positions are margined
against a mark price, and a position whose margin ratio falls below maintenance is force-closed
by the protocol. The resulting trades are recorded in the node's fill stream with a
\texttt{liquidation} field carrying the liquidated user, the mark price and the liquidation
method, alongside the fill's own price, size, starting position and realised profit and loss.

Two properties make the archive suitable for this measurement. It is \textbf{complete}: every
fill is present, with no publication rate limit, because the data is the chain's own record
rather than a broadcast feed. And it is \textbf{attributed}: each fill names the liquidated
account, so fills belonging to one forced close can be grouped without inference.

The archive is hosted as requester-pays object storage, partitioned by hour.

\subsection{Sampling design}

A full census of the archive year is approximately 8{,}760 hours and on the order of 300\,GB of
egress. We instead sample \textbf{four fixed hours per day---02, 08, 14 and 20 UTC---across all
365 archive days}, which is uniform in time and inside the free egress allowance.

The hours are \emph{fixed rather than chosen}. This matters more than the sample size. An earlier
version of this measurement used twelve hours, two of which had been selected because they
contained known liquidation cascades, and therefore over-represented large episodes in a study
whose subject is the size dependence of a bias. Fixing the hours removes that selection directly
rather than correcting for it afterwards.

The sample comprises \textbf{351{,}540 episodes} reconstructed from \textbf{2{,}010{,}042 fills},
across \textbf{380 instruments} and \textbf{151{,}730 distinct accounts}, totalling
\textbf{\$15.53 billion} of liquidated notional.

\subsection{Instrument segments}

Hyperliquid hosts both venue-listed perpetuals (``majors'') and builder-deployed markets under
its HIP-3 mechanism, identified by a \texttt{dex:} prefix in the instrument name. HIP-3 markets
are thinner and separately parameterised, so we report them separately throughout: they carry
17.7\% of episodes but 9.3\% of notional---more liquidations, smaller.

\section{What counts as one liquidation}
\label{sec:unit}

The natural unit is the \emph{episode}: one account's forced close of one position. The archive
does not label episodes, so they must be reconstructed, and the reconstruction rule is a
researcher degree of freedom. We therefore do not defend a single rule. We count under six, and
report the spread.

\begin{table}[htbp]
\centering
\caption{Count inflation is insensitive to the reconstruction rule. Twelve-hour sample,
24{,}566 fills.}
\label{tab:units}
\begin{tabular}{lrr}
\toprule
unit & count & inflation \\
\midrule
fills & 24{,}566 & 1.0$\times$ \\
(account, transaction) & 6{,}412 & 3.8$\times$ \\
episode, 1\,s gap & 6{,}582 & 3.7$\times$ \\
episode, 5\,s gap & 6{,}546 & 3.8$\times$ \\
episode, 60\,s gap & 6{,}533 & 3.8$\times$ \\
(account, instrument, hour) & 6{,}465 & 3.8$\times$ \\
\bottomrule
\end{tabular}
\end{table}

Grouping by transaction, by a one-second gap, or by the whole hour lands within 2\% of the same
answer (Table~\ref{tab:units}). Tranching happens \emph{inside} a transaction, not across
transactions separated in time, so the choice of rule is not doing the work. We adopt
(account, transaction) for the remainder.

One caution earned in the course of this work: a transaction hash is not a liquidation. A single
transaction can carry forced closes for several accounts, so grouping by hash alone
under-counts. The grouping key must include the account.

\section{Count inflation}
\label{sec:count}

On the one-year stratified sample, counting fills instead of episodes inflates the event count
by a factor of

\[
\frac{2{,}010{,}042}{351{,}540} = \mathbf{5.72}.
\]

The factor is comparable across segments: 5.78$\times$ for majors (289{,}283 episodes from
1{,}672{,}034 fills) and 5.43$\times$ for HIP-3 (62{,}257 from 338{,}008).

\subsection{The inflation is size-dependent}

A constant multiplier would be a nuisance parameter. This one is not constant.

\begin{table}[htbp]
\centering
\caption{Tranching grows with episode size. One-year sample, 351{,}540 episodes.}
\label{tab:size}
\begin{tabular}{lrrr}
\toprule
size bucket & episodes & median fills & mean fills \\
\midrule
p0--50 \ (\$0 -- \$1{,}117) & 175{,}770 & 2 & 2.19 \\
p50--90 \ (\$1{,}117 -- \$39{,}307) & 140{,}616 & 2 & 3.94 \\
p90--99 \ (\$39{,}307 -- \$549{,}060) & 31{,}638 & 8 & 19.11 \\
\textbf{p99--100} \ (\$549{,}060 -- \$194{,}115{,}094) & \textbf{3{,}516} & \textbf{72} & \textbf{132.27} \\
\bottomrule
\end{tabular}
\end{table}

The rank correlation between log notional and log fills per episode is $+0.545$. A typical
liquidation is two fills; the largest are a median of 72 and a mean of 132, with a maximum of
4{,}776.

The consequence is stated most sharply as a share: \textbf{the top 1\% of episodes generate
23.1\% of all fills}. A dataset of fills is therefore not a dataset of liquidations with a
scaling factor---it is a dataset in which the largest events have been replicated most.

\subsection{Concentration}

Liquidated notional is concentrated to a degree that compounds the above: the top 1\% of
episodes carry \textbf{67.3\%} of notional, the top 5\% carry 85.6\%, and the top 10\% carry
92.6\%. The events that dominate the economics are the events the fill-count distorts most.

\section{Size compression}
\label{sec:compression}

Fill-counting also changes the measured size distribution, because an episode of a given size
appears as several smaller records.

\begin{table}[htbp]
\centering
\caption{Fill-counting compresses the size distribution. Measured with per-fill notionals on the
twelve-hour sample; see Section~\ref{sec:limits}.}
\label{tab:compression}
\begin{tabular}{lrrr}
\toprule
quantile & true episode & observed fill & factor \\
\midrule
p50 & \$661 & \$424 & 1.6$\times$ \\
p90 & \$18{,}259 & \$8{,}842 & 2.1$\times$ \\
p99 & \$214{,}474 & \$63{,}812 & \textbf{3.4$\times$} \\
p99.9 & \$1{,}969{,}302 & \$180{,}161 & \textbf{10.9$\times$} \\
maximum & \$3{,}838{,}207 & \$899{,}224 & 4.3$\times$ \\
\bottomrule
\end{tabular}
\end{table}

The compression grows through the distribution, reaching an order of magnitude at the 99.9th
percentile (Table~\ref{tab:compression}). In this sample the largest single liquidation of
\$3.84M appears in the fill record as pieces of at most \$899k.

At one-year scale the largest episode is \textbf{\$194{,}115{,}094}, decomposed into 4{,}776
fills. Per-fill notionals were not retained at year scale, so Table~\ref{tab:compression} is
reported from the twelve-hour sample; this is the principal data limitation of the paper and is
straightforward to remove.

\section{The episode-size distribution}
\label{sec:tail}

\subsection{What holds}

The tail is unambiguously heavy. Against a Pareto reference fitted on the same exceedances, an
exponential specification is rejected decisively (Vuong $R = +11.96$, $p < 10^{-4}$, in favour of
Pareto). But the Pareto fit is itself rejected: a Kolmogorov--Smirnov test with
parametric-bootstrap $p$-values rejects it at every threshold tried, including at the threshold
selected by the Clauset--Shalizi--Newman procedure ($p = 0.010$).

Both two-parameter alternatives beat it. Fitting conditional on $x \ge x_{\min}$ and comparing by
Vuong's non-nested likelihood-ratio test, lognormal and Weibull each dominate Pareto at
\textbf{all 14 estimable thresholds}, with $x_{\min}$ ranging from \$14{,}912 to \$8{,}845{,}704
---nearly three orders of magnitude. An exponential never wins at any threshold.

\subsection{What does not: the tail cannot be named}

Selecting $x_{\min}$ by minimising the KS distance of the power-law fit gives \$560{,}627
($k = 3{,}104$, $\hat\alpha = 0.960$). A nonparametric bootstrap over 500 resamples puts a 95\%
interval on that selection of $[\$193{,}877,\ \$986{,}963]$---\textbf{a factor of five}. The
threshold is barely identified by the data, which is why we report the ranking across a grid of
thresholds rather than at a single selected one.

Across that grid the ranking between the two survivors \emph{reverses}:

\begin{table}[htbp]
\centering
\caption{Which alternative wins depends on where the tail is cut. Vuong $R$ for lognormal
against Weibull; $R<0$ favours Weibull.}
\label{tab:reverse}
\begin{tabular}{lrl}
\toprule
$x_{\min}$ range & Vuong $R$ & verdict \\
\midrule
\$199k -- \$1.33M & $-2.74$ to $-3.67$ \ ($p<0.01$) & Weibull \\
\$1.85M -- \$8.85M & $-1.00$ to $-1.86$ \ ($p>0.05$) & indistinguishable \\
\$14.9k -- \$28.3k & $+13.42$ to $+15.92$ \ ($p \ll 0.001$) & lognormal \\
\bottomrule
\end{tabular}
\end{table}

Weibull wins in the far tail, lognormal wins decisively deeper in the body, and the crossover
falls inside a band where neither is identified. \textbf{Any name given to this tail would
report the analyst's threshold rather than the distribution.}

We take this to be the substantive finding of the section rather than a gap in it.

\subsection{A band where nothing fits}

Over $x_{\min} \in [\$49\mathrm{k},\ \$163\mathrm{k}]$, both alternatives converge on a
parameter bound---lognormal with $\mu \to -\infty$, Weibull with $\lambda \to 0$. Both limits are
the power-law limit of their own family. In that band the extra parameter buys nothing and the
models are unidentified, while the power law they are collapsing onto is itself rejected by the
KS test.

So over that band \emph{no candidate here describes the data}. A model ranking cannot report
this, because a ranking always returns a winner. We flag it because a reader who takes the
winning family from a comparison table has no way to see it.

\subsection{Withdrawn estimates}

Three tail-index estimates were produced in the course of this work---$\alpha \approx 1.15$,
then $\alpha \in [0.9, 1.2]$, then a Hill-plot plateau at $\alpha \approx 0.93$---and all three
are withdrawn. The Hill estimator estimates a Pareto exponent, and the tail is not Pareto; drift
of the Hill index with tail depth, initially read as a measurement defect, is the known
signature of a lognormal-like tail. Consequences drawn from those numbers, including ``infinite
variance'' and ``the sample mean is not a stable statistic,'' are withdrawn with them: both are
properties of a fitted Pareto with $\alpha < 2$, and that fit is rejected.

The count and compression results in Sections~\ref{sec:count}--\ref{sec:compression} are
unaffected. They compare two curves computed identically on the same data and assume no
parametric form.

\section{Why this does not transfer to centralised-exchange feeds}
\label{sec:cex}

The bias measured here is specific to a \emph{fill}-level record. Centralised venues publish
liquidation \emph{orders}, so there is no tranching to inflate. Their feeds carry a different
defect.

\begin{table}[htbp]
\centering
\caption{What the major centralised feeds publish, per venue documentation.}
\label{tab:cex}
\begin{tabular}{llll}
\toprule
venue & stream & unit & limitation \\
\midrule
Binance & \texttt{forceOrder} & order & only the \emph{largest} order per 1000\,ms per symbol \\
OKX & \texttt{liquidation-orders} & order & at most one update per second per contract \\
Bybit (legacy) & \texttt{liquidation} & order & at most 1/s per symbol; deprecated \\
Bybit (current) & \texttt{allLiquidation} & order & complete, pushed every 500\,ms \\
\bottomrule
\end{tabular}
\end{table}

Binance does not publish a sample of liquidations. It publishes a sequence of \emph{per-second
maxima}, which is a different statistical object with three consequences:

\begin{itemize}
\item \textbf{event counts are unrecoverable}---nothing in the feed states how many liquidations
      a given snapshot suppressed;
\item \textbf{the body of the distribution is destroyed}---small liquidations never appear at
      all;
\item \textbf{the tail index may survive}---block maxima are a valid extreme-value approach, and
      a block maximum inherits the tail index of its parent distribution.
\end{itemize}

That last point matters for interpreting existing results: a study of Binance liquidation data
can produce a defensible tail estimate and a worthless event count from the same feed. The
distinction is between statistics that depend on counting or on the body, which are compromised,
and tail-shape statistics, which may not be.

We note this from venue documentation, not from a re-analysis of published work, and make no
claim about any specific paper.

\section{Implications}

\paragraph{For measurement.} Any statistic over crypto liquidation data should state its unit.
On a fill-level record the correction is a factor of 5.72 on counts, but because the factor is
size-dependent it cannot be applied as a scalar: the correct procedure is to reconstruct
episodes, which requires an attributed record.

\paragraph{For cascade studies.} Cascade detection typically thresholds on event counts or
event rates within a window. Both are inflated most during the large episodes that constitute
cascades, so a fill-based cascade measure partly detects its own tranching.

\paragraph{For risk models.} A size distribution fitted on fills is compressed by 3.4$\times$ at
the 99th percentile and 10.9$\times$ at the 99.9th. A model calibrated on that distribution
understates the size of the largest forced closes it is meant to survive.

\paragraph{For the tail literature.} We would encourage reporting the model ranking across a
range of thresholds rather than at a selected one, and reporting where no candidate fits. In this
dataset a single-threshold analysis would have supported naming the tail; the grid shows the
name would have been the threshold's.

\section{Limitations}
\label{sec:limits}

\begin{itemize}
\item \textbf{Compression is measured on twelve hours}, not on the year, because per-fill
      notionals were not retained at year scale. The count and tranching results are year-scale.
      Re-collection is straightforward.
\item \textbf{The sample is stratified, not exhaustive}: four fixed hours per day. Uniform in
      time, but a cascade falling entirely outside those hours is invisible.
\item \textbf{One venue, one year.} Whether the inflation factor is stable over time, or trends
      with venue growth, is untested.
\item \textbf{The goodness-of-fit bootstrap fixes $x_{\min}$} inside each synthetic sample rather
      than re-selecting it, which makes the reported $p$ conservative toward rejection.
\item \textbf{Four candidate families.} A fifth---generalised Pareto, log-gamma---might describe
      the band where none of ours does. None was tried.
\item \textbf{The literature review is incomplete.} We position this work against venue
      documentation and against the statistical methodology it uses, but have not surveyed the
      empirical crypto-liquidation literature systematically. Claims of novelty should be read
      as claims about what we measured, not about what has been measured before.
\end{itemize}

\section{Reproducibility}
\label{sec:repro}

All code, reduced datasets and experimental records are public at
\url{https://github.com/QuasareumGroup/hyperliquid-microstructure}.

Each experiment is pre-registered: hypothesis, predictions and falsification criteria are
committed \emph{before} the run, and the commit ordering is checkable in the repository history.
Corrections and withdrawals remain in the file that made the original claim rather than being
edited away; \texttt{experiments/FINDINGS.md} carries the current state of every claim and the
reason for each of the fourteen that have been withdrawn or requalified.

Published datasets carry a derived account identifier---a truncated SHA-256 of the address---in
place of the address itself. Addresses are public on-chain, so this is not confidentiality; it
avoids publishing a compiled map of 151{,}730 accounts to their losses. Every result reproduces
bit for bit on the derived identifiers. The hash is unsalted and deterministic, so a known
address can still be tested for membership: these datasets are pseudonymous, not anonymised.

\begin{thebibliography}{9}

\bibitem{clauset2009}
A.~Clauset, C.~R. Shalizi, and M.~E.~J. Newman.
\newblock Power-law distributions in empirical data.
\newblock \emph{SIAM Review}, 51(4):661--703, 2009.

\bibitem{vuong1989}
Q.~H. Vuong.
\newblock Likelihood ratio tests for model selection and non-nested hypotheses.
\newblock \emph{Econometrica}, 57(2):307--333, 1989.

\bibitem{hill1975}
B.~M. Hill.
\newblock A simple general approach to inference about the tail of a distribution.
\newblock \emph{Annals of Statistics}, 3(5):1163--1174, 1975.

\bibitem{hyperliquid}
Hyperliquid.
\newblock Documentation: node data, liquidations, and the funding mechanism.
\newblock \url{https://hyperliquid.gitbook.io/hyperliquid-docs}.

\bibitem{binance}
Binance.
\newblock USD-M futures WebSocket streams: liquidation order streams.
\newblock \url{https://developers.binance.com/docs/derivatives}.

\end{thebibliography}

\end{document}
